\documentclass{article}
\bibliographystyle{plainnat}
\usepackage{amsmath, amssymb}
\usepackage{url}
\usepackage{natbib}
\usepackage{minted}   % For code highlighting
\usepackage{graphicx}
\usepackage{enumitem}
\usepackage{titlesec}
\usepackage[margin=1in]{geometry}

% ============================================================
% Title and Author – adapt as needed
% ============================================================
\title{Your Title Here}
\author{
  {\Large Your Name} \\
  \vspace{0.2em}
  {\small Your Affiliation, Institution} \\
  {\small \texttt{your.email@domain}}
}
\date{\today}

\begin{document}

\maketitle

% ----------------------------------------------------------------------
% Optional abstract / introduction
% ----------------------------------------------------------------------
\section*{Introduction}
Write your introduction here. This template retains the exact formatting 
(margins, fonts, packages) from the original GRPO paper.

% ----------------------------------------------------------------------
% Main sections – adapt structure as needed
% ----------------------------------------------------------------------
\section{First Section}
\label{sec:first}

Here you can include theoretical details, equations, and figures.

\begin{figure}[h]
    \centering
    \includegraphics[width=0.8\linewidth]{example-image} % replace with your image
    \caption{Example figure caption.}
    \label{fig:example}
\end{figure}

\subsection{A Subsection}
You can use \texttt{minted} for code snippets:

\begin{minted}[fontsize=\small, bgcolor=gray!5, frame=lines, framesep=2mm, baselinestretch=1.1, linenos]{python}
def hello():
    print("Hello, world!")
\end{minted}

Equations are supported:
\[
    E = mc^2
\]

\section{Second Section}
\label{sec:second}

More content here.

\subsection{Key Formulas}
\begin{enumerate}
    \item Advantage: \( A_i = \frac{r_i - \bar{r}}{\sigma_r + \epsilon} \)
    \item Probability ratio: \( \text{ratio}_{i,t} = \frac{\pi_\theta(o_{i,t} \mid q, o_{i,<t})}{\pi_{\theta_{\text{old}}}(o_{i,t} \mid q, o_{i,<t})} \)
\end{enumerate}

\section{Conclusion}
Summarize your work.

% ----------------------------------------------------------------------
% Bibliography – keep the same style and use \bibitem or .bib file
% ----------------------------------------------------------------------
\begin{thebibliography}{9}

\bibitem{key1}
Author. \textit{Title}. Available at: \url{https://example.com}. Accessed: Date.

\bibitem{key2}
Another Author. \textit{Book or Paper}. Journal, Year.

\end{thebibliography}

\end{document}
